% GENERATED FILE. Edit canonical lessons, then run npm run generate:books.
% canonical-source-hash: fnv1a64:897e4b549d219384
% canonical-lessons: ZH-C509-dangjia, ZH-C509-fumu, ZH-W509-mu-mother, ZH-C509-munu, ZH-C509-muyu, ZH-C509-houma

\chapter{Parents}
\label{ch:zh-parents}
\hlchaptermodality{509}

\begin{chapteropening}
\textbf{By the end of this chapter:} \emph{I can talk about who runs the household, my parents, mothers and daughters, my mother tongue and stepmothers.}

Complete the last lesson of chapter 509: I can talk about who runs the household, my parents, mothers and daughters, my mother tongue and stepmothers.
\end{chapteropening}

\section[\texorpdfstring{\zh{当家} (dāngjiā)}{当家 (dāngjiā)}]{\zh{当家} (dāngjiā) — to run the household}

\label{lesson:ZH-C509-dangjia}



\begin{quote}
Before the new one: say the Mandarin for to pick on (someone weaker), then the Mandarin for sensible (of a child).
\end{quote}

\subsection*{You'll want to know: \zh{当家}}
\textbf{\zh{当家}} — \emph{dāngjiā} — \textquotedblleft{}to run the household\textquotedblright{}.

\begin{quote}
\textbf{\zh{我们家是我妈妈当家}} \emph{wǒmen jiā shì wǒ māma dāngjiā} — In our family, it's my mum who runs things.
\end{quote}

\begin{cousinweb}
\textbf{\zh{当家}} (\emph{dāngjiā}) is built from \zh{当} 'to be in charge' and \zh{家} 'a home', both inherited from Middle Chinese, \emph{tang} and \emph{kae} (Guangyun, in Baxter's transcription), which glosses the first 'to be in charge' and the second 'a dwelling'.
\end{cousinweb}

\begin{grammarlens}[title={What you've built}]
One more word for saying how it was.
\end{grammarlens}

\subsection*{Your turn}
\begin{itemize}
\raggedright
  \item \emph{Say it:} \emph{dāngjiā}
  \item \emph{Say it:} \emph{dāngjiā}, once more
  \item \emph{Say it:} \emph{wǒmen jiā shì wǒ māma dāngjiā} — In our family, it's my mum who runs things.
  \item \emph{Recall:} say the Mandarin for to pick on (someone weaker), then the Mandarin for sensible (of a child), then say \emph{dāngjiā} again
\end{itemize}

\begin{tcolorbox}[breakable,colback=teal!4,colframe=teal!35!black,title={Before you move on}]
What does \zh{当家} mean? (\textquotedblleft{}To run the household\textquotedblright{}.) Say it once more.
\end{tcolorbox}
\section[\texorpdfstring{\zh{父母} (fùmǔ)}{父母 (fùmǔ)}]{\zh{父母} (fùmǔ) — parents; one's mother and father}

\label{lesson:ZH-C509-fumu}



\begin{quote}
Before the new one: say the Mandarin for sensible (of a child), then the Mandarin for to run the household.
\end{quote}

\subsection*{You'll want to know: \zh{父母}}
\textbf{\zh{父母}} — \emph{fùmǔ} — \textquotedblleft{}parents; one's mother and father\textquotedblright{}.

\begin{quote}
\emph{wǒ gēn fùmǔ zhù le èrshí nián} — I lived with my parents for twenty years.
\end{quote}

The new character, in pieces: \zh{母} — five strokes: a frame of two, one down the left that turns right along the bottom and one along the top that turns down the right side and hooks; a dot inside the top half; a long bar across; a dot inside the bottom half. The next lesson puts it in your hand.

\begin{cousinweb}
\textbf{\zh{父母}} (\emph{fùmǔ}) is built from \zh{父} 'a father' and \zh{母} 'a mother', both inherited from Middle Chinese, \emph{bjuX} and \emph{muwX} (Guangyun, in Baxter's transcription), which glosses the second 'a mother' and, after an older book, reads its two dots as a picture of a nursing mother's breasts.
\end{cousinweb}

\begin{grammarlens}[title={What you've built}]
One more word for saying how it was.
\end{grammarlens}

\subsection*{Your turn}
\begin{itemize}
\raggedright
  \item \emph{Say it:} \emph{fùmǔ}
  \item \emph{Say it:} \emph{fùmǔ}, once more
  \item \emph{Say it:} \emph{wǒ gēn fùmǔ zhù le èrshí nián} — I lived with my parents for twenty years.
  \item \emph{Recall:} say the Mandarin for sensible (of a child), then the Mandarin for to run the household, then say \emph{fùmǔ} again
\end{itemize}

\begin{tcolorbox}[breakable,colback=teal!4,colframe=teal!35!black,title={Before you move on}]
What does \zh{父母} mean? (\textquotedblleft{}Parents; one's mother and father\textquotedblright{}.) Say it once more.
\end{tcolorbox}
\section[\texorpdfstring{\zh{母} (mǔ)}{母 (mǔ)}]{\zh{母} — a frame, two dots and a long bar}

\label{lesson:ZH-W509-mu-mother}



\begin{quote}
The word first, then the character inside it.

\begin{itemize}
\raggedright
  \item \emph{Recall:} say \emph{fùmǔ} — \textbf{R1}, one lesson back, before it decays
\end{itemize}

\begin{quote}
\zh{父} \zh{女}
\end{quote}

\emph{Read it:} each one aloud, with nothing said first — \emph{fù}, \emph{nǚ}
\end{quote}

\subsection*{Script — a frame, two dots and a long bar}
\hlblockfigure{\detokenize{figures/ZH-W509-mu-mother-filmstrip.pdf}}{How \zh{母} is written, stroke by stroke}

\begin{quote}
\zh{母}
\end{quote}

\textbf{The frame, two strokes.} One that runs down the left side and turns right along the bottom. Then one that runs along the top, turns down the right side and hooks to the left at the foot.

\textbf{A dot} inside the top half.

\textbf{A long bar} right across the middle, wider than the frame on both sides.

\textbf{A dot} inside the bottom half.

2 + 1 + 1 + 1 = five.

\subsection*{Your turn}
\begin{itemize}
\raggedright
  \item \emph{Write it:} \zh{母}, all five, with the model in front of you
  \item \emph{Cover:} write it again from memory, and count the pieces, not the strokes
\end{itemize}

\begin{tcolorbox}[breakable,colback=teal!4,colframe=teal!35!black,title={Before you move on}]
How many strokes does \zh{母} have? (\textbf{Five}.) How many dots? (\textbf{Two, one above the long bar and one below}.) Which stroke is widest? (\textbf{The long bar across the middle}.)

Source: \href{https://raw.githubusercontent.com/chanind/hanzi-writer-data/68d10a4b21150cae5e1ebbd223eed289cf32d90c/data/\%E6\%AF\%8D.json}{Hanzi Writer Data \zh{母}, pinned stroke paths and medians}.
\end{tcolorbox}
\section[\texorpdfstring{\zh{母女} (mǔnǚ)}{母女 (mǔnǚ)}]{\zh{母女} (mǔnǚ) — mother and daughter}

\label{lesson:ZH-C509-munu}



\begin{quote}
Before the new one: say the Mandarin for to run the household, then the Mandarin for one's mother and father.
\end{quote}

\subsection*{You'll want to know: \zh{母女}}
\textbf{\zh{母女}} — \emph{mǔnǚ} — \textquotedblleft{}mother and daughter\textquotedblright{}.

\begin{quote}
\textbf{\zh{母女两个人都喜欢去舞会}} \emph{mǔnǚ liǎng ge rén dōu xǐhuan qù wǔhuì} — Mother and daughter both love going to dances.
\end{quote}

Two third tones meet in this word: mǔ is said with a rising tone before nǚ, and the spelling keeps both marks.

\begin{cousinweb}
\textbf{\zh{母女}} (\emph{mǔnǚ}) is built from \zh{母} 'a mother' and \zh{女} 'a daughter', both inherited from Middle Chinese, \emph{muwX} and \emph{nrjoX} (Guangyun, in Baxter's transcription).
\end{cousinweb}

\begin{grammarlens}[title={What you've built}]
One more word for saying how it was.
\end{grammarlens}

\subsection*{Your turn}
\begin{itemize}
\raggedright
  \item \emph{Say it:} \emph{mǔnǚ}
  \item \emph{Say it:} \emph{mǔnǚ}, once more
  \item \emph{Say it:} \emph{mǔnǚ liǎng ge rén dōu xǐhuan qù wǔhuì} — Mother and daughter both love going to dances.
  \item \emph{Recall:} say the Mandarin for to run the household, then the Mandarin for one's mother and father, then say \emph{mǔnǚ} again
\end{itemize}

\begin{tcolorbox}[breakable,colback=teal!4,colframe=teal!35!black,title={Before you move on}]
What does \zh{母女} mean? (\textquotedblleft{}Mother and daughter\textquotedblright{}.) Say it once more.
\end{tcolorbox}
\section[\texorpdfstring{\zh{母语} (mǔyǔ)}{母语 (mǔyǔ)}]{\zh{母语} (mǔyǔ) — mother tongue}

\label{lesson:ZH-C509-muyu}



\begin{quote}
Before the new one: say the Mandarin for one's mother and father, then the Mandarin for a mother and daughter.
\end{quote}

\subsection*{You'll want to know: \zh{母语}}
\textbf{\zh{母语}} — \emph{mǔyǔ} — \textquotedblleft{}mother tongue\textquotedblright{}.

\begin{quote}
\emph{wǒ de mǔyǔ shì Yīngyǔ} — My mother tongue is English.
\end{quote}

Two third tones meet in this word: mǔ is said with a rising tone before yǔ, and the spelling keeps both marks.

\begin{cousinweb}
\textbf{\zh{母语}} (\emph{mǔyǔ}) is built from \zh{母} 'a mother' and \zh{语} 'speech', both inherited from Middle Chinese, \emph{muwX} and \emph{ngjoX} (Guangyun, in Baxter's transcription).
\end{cousinweb}

\begin{grammarlens}[title={What you've built}]
One more word for saying how it was.
\end{grammarlens}

\subsection*{Your turn}
\begin{itemize}
\raggedright
  \item \emph{Say it:} \emph{mǔyǔ}
  \item \emph{Say it:} \emph{mǔyǔ}, once more
  \item \emph{Say it:} \emph{wǒ de mǔyǔ shì Yīngyǔ} — My mother tongue is English.
  \item \emph{Recall:} say the Mandarin for one's mother and father, then the Mandarin for a mother and daughter, then say \emph{mǔyǔ} again
\end{itemize}

\begin{tcolorbox}[breakable,colback=teal!4,colframe=teal!35!black,title={Before you move on}]
What does \zh{母语} mean? (\textquotedblleft{}Mother tongue\textquotedblright{}.) Say it once more.
\end{tcolorbox}
\section[\texorpdfstring{\zh{后妈} (hòumā)}{后妈 (hòumā)}]{\zh{后妈} (hòumā) — stepmother}

\label{lesson:ZH-C509-houma}



\begin{quote}
Before the new one: say the Mandarin for a mother and daughter, then the Mandarin for one's mother tongue.
\end{quote}

\subsection*{You'll want to know: \zh{后妈}}
\textbf{\zh{后妈}} — \emph{hòumā} — \textquotedblleft{}stepmother\textquotedblright{}.

\begin{quote}
\emph{tā xiǎoshíhou gēn hòumā zhù le wǔ nián} — As a child he lived with his stepmother for five years.
\end{quote}

\begin{cousinweb}
\textbf{\zh{后妈}} (\emph{hòumā}) is built from \zh{后} 'after' and \zh{妈} 'mum', both inherited from Middle Chinese, \emph{huwX} and \emph{muX} (Guangyun, in Baxter's transcription), which glosses the first 'after' and the second 'a mother'. It belongs to everyday speech.
\end{cousinweb}

\begin{grammarlens}[title={What you've built}]
One more word for saying how it was.
\end{grammarlens}

\subsection*{Your turn}
\begin{itemize}
\raggedright
  \item \emph{Say it:} \emph{hòumā}
  \item \emph{Say it:} \emph{hòumā}, once more
  \item \emph{Say it:} \emph{tā xiǎoshíhou gēn hòumā zhù le wǔ nián} — As a child he lived with his stepmother for five years.
  \item \emph{Recall:} say the Mandarin for a mother and daughter, then the Mandarin for one's mother tongue, then say \emph{hòumā} again
\end{itemize}

\begin{tcolorbox}[breakable,colback=teal!4,colframe=teal!35!black,title={Before you move on}]
What does \zh{后妈} mean? (\textquotedblleft{}Stepmother\textquotedblright{}.) Say it once more.
\end{tcolorbox}
